\documentclass[11pt]{article}

\usepackage[margin=1in]{geometry}
\usepackage{amsmath,amssymb}
\usepackage{graphicx}
\usepackage{booktabs}
\usepackage{hyperref}
\usepackage{url}
\usepackage{cite}
\usepackage{authblk}

\title{A Neural Speech Decoding Framework Leveraging Deep Learning and a
Differentiable Speech Synthesizer for ECoG-Based Brain-Computer Interfaces}

\date{}

\begin{document}
\maketitle

\begin{abstract}
Restoring intelligible, natural-sounding speech to individuals with severe
neurological deficits remains one of the defining goals of modern
brain-computer interface (BCI) research. Progress in this area is constrained
by three recurring obstacles: the small amount of paired neural and speech
data that can be collected from any single participant, the widespread
reliance on non-causal operations that tacitly exploit auditory feedback
unavailable at deployment, and the scarcity of openly released pipelines
that permit head-to-head comparison of decoder architectures. We address
these gaps with a two-stage neural speech decoding framework that couples
an ECoG Decoder, implemented with interchangeable 3D ResNet, 3D Swin
Transformer, or LSTM backbones, to a fully differentiable Speech Synthesizer
defined over a compact space of $18$ interpretable acoustic parameters.
A companion speech-to-speech auto-encoder is first trained on speech alone
to produce reference parameter trajectories that then supervise ECoG Decoder
training, allowing spectral reconstruction loss to be back-propagated
end-to-end. Applied to $N=48$ neurosurgical participants, the causal
ResNet decoder achieved a mean spectrogram Pearson correlation coefficient
(PCC) of $0.798$ against the ground-truth spectrogram, closely followed by
the causal Swin decoder at PCC $0.796$; neither causal variant differed
significantly from its non-causal counterpart (Wilcoxon signed-rank
$p=0.1022$ and $p=0.2794$), while the LSTM showed a significant causal
deficit ($p=0.0007$). Right-hemisphere participants ($n=16$) reached PCC
$0.790$ (ResNet) and $0.781$ (Swin), statistically indistinguishable from
left-hemisphere results. Together, these findings establish a reproducible,
causal, natural-sounding ECoG-to-speech pipeline, provide the first
systematic evidence for right-hemisphere speech prostheses, and release the
full framework as open-source software to accelerate community progress.
\end{abstract}

\section{Introduction}

Loss of spoken language is among the most disabling consequences of
neurological disease, profoundly limiting social participation and
employment for survivors of stroke, amyotrophic lateral sclerosis, and
severe brain injury. Brain-computer interface systems offer a principled
path toward restoring communication by translating cortical activity back
into intelligible speech output, and subdural electrocorticography (ECoG)
recorded intraoperatively from epilepsy surgery patients has become a
central substrate for this endeavour. ECoG grids provide spatially and
temporally precise measurements of cortical activity, and a growing body
of work has shown that the high-gamma envelope of these recordings carries
phonemic, articulatory, and spectrotemporal correlates of spoken
speech~\cite{pasley2012reconstructing,bouchard2014control,mugler2014direct,herff2015brain,flinker2015redefining,wang2021dynamics}.
Despite this progress, converting cortical activity into natural-sounding
speech at latencies compatible with conversational BCIs remains an open
problem, constrained by the limited amount of paired neural and speech
data that can be collected from a single participant and by the high
trial-to-trial variability of speech production.

Early ECoG-to-speech decoders relied on linear regression from high-gamma
activity onto acoustic spectrograms or phonetic features. Such models
typically reached a Pearson correlation coefficient of approximately $0.6$
between original and reconstructed spectrograms, offering interpretability
and modest data requirements at the cost of audible intelligibility. More
recent deep learning approaches have improved on this baseline with
convolutional, recurrent, and hybrid
architectures~\cite{akbari2019towards,anumanchipalli2019speech,makin2020machine}.
These systems differ both in the intermediate latent representation they
use to describe speech and in the naturalness of the audio they produce.
One prominent line of work decodes cortical activity into an articulatory
kinematic space that is subsequently mapped to
waveforms~\cite{anumanchipalli2019speech}; the resulting audio is accurate
but adopts a generic voice unrelated to the participant. A complementary
line uses wavenet-style vocoders~\cite{akbari2019towards}, generative
adversarial networks, or unit selection to generate more natural audio,
often trading off decoding accuracy. In parallel, clinical speech BCIs
have delivered striking demonstrations in individuals with severe paralysis:
intracortical microelectrode arrays now support high-vocabulary text
decoding at conversational
speeds~\cite{willett2023highperformance}, and high-density ECoG systems
have been used to restore synchronous text, audio, and avatar control in
anarthric participants~\cite{moses2021neuroprosthesis,metzger2023highperformance,silva2024bilingual}.
These achievements confirm the clinical feasibility of speech prostheses
but typically depend on non-causal temporal processing, which implicitly
exploits post-utterance auditory feedback unavailable at deployment.

A second gap concerns reproducibility. Most state-of-the-art ECoG speech
decoding pipelines are not released in a form that supports independent
replication, and direct head-to-head comparison across architectures,
causality modes, electrode densities, or hemispheres is therefore
difficult. Nearly all prior work has additionally concentrated on
left-hemisphere, language-dominant cortex, even though a sizeable fraction
of patients who would benefit from a speech BCI have sustained
left-hemisphere damage and might be served instead by a right-hemisphere
prosthesis~\cite{lorach2023walking}. The absence of multi-participant,
causal, open-source benchmarks leaves the field without a common reference
point against which new decoders can be measured.

To address these issues, we propose a two-stage ECoG-to-speech framework
organised around two components. The first component is an ECoG Decoder
that maps cortical high-gamma activity to a compact set of $18$
interpretable speech parameters per time step, including pitch, voicing,
formant frequencies and amplitudes, loudness, and broadband noise
descriptors. The second component is a fully differentiable Speech
Synthesizer that maps those parameters back to spectrograms. A companion
speech-to-speech auto-encoder, built from a Speech Encoder and the same
Speech Synthesizer, is pre-trained on participant speech alone to generate
reference parameter trajectories that supervise ECoG Decoder training.
Because the Speech Synthesizer is differentiable, spectral reconstruction
loss can be back-propagated end-to-end, and the combination of a
low-dimensional parameter space with reference supervision mitigates data
scarcity. Within this framework we provide three interchangeable ECoG
Decoder backbones -- 3D ResNet, 3D Swin Transformer, and LSTM -- each
usable in causal or non-causal temporal modes.

Our contributions are as follows. We introduce the two-stage training
pipeline and its differentiable Speech Synthesizer, and demonstrate across
$N=48$ ECoG participants that the ResNet backbone achieves the highest
spectrogram PCC ($0.804$ non-causal, $0.798$ causal) with Swin a close
second ($0.793$ non-causal, $0.796$ causal). We show that causal ResNet
and causal Swin models are statistically indistinguishable from their
non-causal counterparts (Wilcoxon signed-rank $p=0.1022$ and $p=0.2794$),
whereas causal LSTM is significantly worse than non-causal LSTM
($p=0.0007$). We present the first systematic evaluation of
right-hemisphere decoding with modern deep architectures, reporting PCCs
of $0.790$ (ResNet) and $0.781$ (Swin) that do not differ significantly
from left-hemisphere results. We further show that low-density clinical
ECoG grids support decoding performance comparable to hybrid high-density
grids, and provide an occlusion-based contribution analysis that
distinguishes the cortical substrates of causal and non-causal decoders.
Finally, we release the complete pipeline as open-source software to
facilitate community-wide replication and comparison.

\section{Related Work}

\paragraph{Linear ECoG-to-speech decoding and cortical speech representation.}
Classical ECoG-to-speech studies use linear regression from high-gamma
activity to acoustic spectrograms or phonetic
features~\cite{pasley2012reconstructing,mugler2014direct,herff2015brain}.
These models typically reach a PCC near $0.6$ between original and
reconstructed spectrograms while remaining easy to interpret and
train. Complementary work has mapped the cortical representation of
speech across peri-sylvian regions, showing that ventral sensorimotor
cortex encodes spoken vowel
acoustics~\cite{bouchard2014control}, that Broca's area coordinates
articulation-related cortical networks~\cite{flinker2015redefining,wang2021dynamics},
and that auditory cortex predictively tracks speech
envelopes~\cite{golumbic2013visual,yildiz2016predictive}. Together these
findings identify superior temporal, inferior frontal, pre-central, and
post-central gyri as prime candidate regions for deep ECoG decoders.

\paragraph{Deep-learning speech BCIs.}
Subsequent work has adopted deep neural networks. Anumanchipalli et
al.~\cite{anumanchipalli2019speech} introduced a recurrent pipeline that
decodes cortical activity into an articulatory kinematic space and
synthesises speech from it; the latent space supports partial
cross-participant transfer but yields a non-naturalistic voice. Akbari et
al.~\cite{akbari2019towards} used a deep network to reconstruct speech
from human auditory cortex and improved digit-level intelligibility by
$65\%$ over a linear baseline. Makin et
al.~\cite{makin2020machine} decoded ECoG directly to text with
word-error rates as low as $3\%$ on a closed sentence set. Clinical
systems have extended these ideas to participants with severe paralysis:
Moses et al.~\cite{moses2021neuroprosthesis} delivered real-time word and
sentence decoding in an anarthric individual; Willett et
al.~\cite{willett2023highperformance} achieved $62$ words per minute and a
$9.1\%$ word-error rate on a $50$-word vocabulary using intracortical
arrays; and Metzger et al.~\cite{metzger2023highperformance} demonstrated
synchronous text, audio, and avatar output from high-density ECoG. Silva
et al.~\cite{silva2024bilingual} further showed that a shared articulatory
representation supports bilingual English--Spanish decoding. These
demonstrations confirm clinical feasibility but generally report single
participants, rely on non-causal operations at inference, and rarely
release code for independent replication.

\paragraph{Backbone architectures and evaluation metrics.}
Our decoder backbones are grounded in three widely used families. The 3D
ResNet~\cite{he2016deep} serves as the convolutional baseline and treats
stacks of ECoG frames as volumetric inputs. The 3D Swin
Transformer~\cite{liu2021swin,liu2022video} extends shifted-window
self-attention to spatio-temporal tokens using $2\times2\times2$ patches
and patch-merging stages. The LSTM~\cite{hochreiter1997lstm} represents
the recurrent family dominant in earlier speech BCIs. All models are
trained with the Adam optimizer~\cite{kingma2015adam}. Output quality is
assessed with the Pearson correlation coefficient of the reconstructed
spectrogram and with the Short-Time Objective Intelligibility family of
metrics~\cite{taal2011algorithm}, including the envelope-power-spectrum
variant~\cite{relanoiborra2016sepsm}. Reference pitch and formant
trajectories are computed with the Praat analysis
toolkit~\cite{boersma2001praat}.

\paragraph{Neural prostheses beyond speech.}
Our work inherits from a broader body of evidence that cortical signals
can drive clinically useful motor prostheses, from early real-time
arm-reach decoding~\cite{wessberg2000realtime} to recent brain--spine
interfaces that restored natural walking after spinal cord
injury~\cite{lorach2023walking}. These successes argue that carefully
engineered decoders can generalise across effectors, and motivate speech
BCIs that exploit both low-density clinical grids and bilateral cortical
coverage -- the two axes we evaluate here.

\section{Methods}

\subsection{ECoG-to-Speech Decoding Framework}

The decoding framework is organised around two components: an ECoG
Decoder and a Speech Synthesizer. The Speech Synthesizer takes $18$
interpretable speech parameters per time step -- pitch $f_0$, voice weight
$\alpha$, loudness $L$, six formants $f_1,\dots,f_6$ with amplitudes,
broadband filter parameters for unvoiced excitation, and a background
noise descriptor -- and maps them to a spectrogram. Training proceeds in
two stages. In stage one, we pre-train a speech-to-speech auto-encoder
consisting of a Speech Encoder and the Speech Synthesizer on each
participant's speech alone; the resulting reference speech parameters
are stored per frame. In stage two, we train the ECoG Decoder on paired
ECoG and speech recordings using a spectrogram reconstruction loss,
supervised additionally by the reference parameters from stage one.
Figure~\ref{fig:architecture} summarises the pipeline.

\begin{figure}[htbp]
\centering
\includegraphics[width=0.7\textwidth]{figures/fig_architecture.pdf}
\caption{Overview of the two-stage neural speech decoding framework.
Stage 1 (bottom) pre-trains a speech-to-speech auto-encoder consisting of a
Speech Encoder and a fully differentiable Speech Synthesizer on the
participant's speech only, producing reference values for 18 interpretable
speech parameters per frame. Stage 2 (top) trains an ECoG Decoder (3D ResNet,
3D Swin Transformer, or LSTM; causal or non-causal) to predict those speech
parameters from ECoG high-gamma activity sampled at 125~Hz. The same
Speech Synthesizer is reused to map the decoded parameters back to a
spectrogram, enabling end-to-end back-propagation of spectral loss. The
reference speech parameters from stage 1 supervise stage 2 through a
reference-parameter loss.}
\label{fig:architecture}
\end{figure}

\subsection{Differentiable Speech Synthesizer}

The Speech Synthesizer combines two pathways. A voice pathway applies $N=6$
formant filters to a harmonic excitation and an unvoice pathway applies a
single broadband filter to white noise; the two are combined with a soft
mix governed by a voice weight $\alpha_t$ at each time step. Prior
linguistics work~\cite{anumanchipalli2019speech} establishes that $N=2$
formants suffice for intelligibility; we retain $N=6$ for more accurate
synthesis and to preserve participant-specific voice quality. Each filter
is parameterised by $M=80$ prototype parameters, and the relationship
between formant centre frequency and bandwidth is captured by a threshold
frequency $f_\theta$, a slope $a$, and a baseline bandwidth $b_0$ that are
shared across the six formant filters and learned during auto-encoder
training. The frequency range $[0, 8000]\,\mathrm{Hz}$ is sampled with
$K=256$ bins for female participants and $K=512$ for male participants.
The total count of learnable parameters in the synthesizer is therefore
$M(N+1)=560$ prototype parameters, three shared bandwidth parameters, and
$K$ background-noise parameters.

\subsection{ECoG Decoder Architectures}

We implement three decoder backbones, each available in causal and
non-causal variants. The 3D ResNet~\cite{he2016deep} treats a sliding
window of ECoG frames as a volumetric input, down-samples the electrode
dimension to $1\times1$ and the temporal dimension to $T/16$, and then
up-samples with transposed convolutions to recover the original temporal
length. The 3D Swin Transformer~\cite{liu2021swin,liu2022video} treats
each $2\times2\times2$ spatio-temporal patch as a token and performs
$2\times2\times2$ down-sampling at the patch-merging stage between
successive transformer blocks. The LSTM~\cite{hochreiter1997lstm} represents
the recurrent family typical of earlier speech BCIs. Causal variants
replace all forward-looking temporal operations with strictly causal ones,
so that decoding at time $t$ depends only on ECoG samples up to and
including $t$.

\subsection{Data Collection and Preprocessing}

We recruited $N=48$ native English speakers (26 female, 22 male) with
refractory epilepsy at NYU Langone. Of these, $43$ participants received
standard low-density (LD) $8\times8$ grids with $64$ macro contacts at
$10\,\mathrm{mm}$ spacing, and $5$ participants received hybrid-density
(HB) grids in which $64$ additional $1\,\mathrm{mm}$ micro-electrodes were
interleaved between the macro contacts, yielding $128$ electrodes with
$5\,\mathrm{mm}$ macro-to-micro centre spacing. $32$ participants had
left-hemisphere coverage and $16$ had right-hemisphere coverage.
Participants completed five speech tasks (auditory repetition, auditory
naming, sentence completion, word reading or visual reading, and picture
naming) over a set of $50$ target words, producing $400$ trials per
participant with a $500\,\mathrm{ms}$ mean speech duration. ECoG signals
were sampled at $2048\,\mathrm{Hz}$, down-sampled to $512\,\mathrm{Hz}$,
and the Hilbert transform was used to extract the high-gamma
($70$--$150\,\mathrm{Hz}$) envelope, which was further down-sampled to
$125\,\mathrm{Hz}$. A $250\,\mathrm{ms}$ silent baseline preceding each
stimulus provided the reference for $z$-scoring, and artefactual and
interictal electrodes were excluded. Reference pitch $f_0$ and formants
$f_1,\dots,f_4$ were estimated with the Praat
toolkit~\cite{boersma2001praat} and resampled to $125\,\mathrm{Hz}$ to
match the ECoG and spectrogram grids.

\subsection{Training and Loss Functions}

The ECoG Decoder is supervised by three complementary loss terms. A
spectral loss $\mathcal{L}_{\mathrm{spec}}$ penalises differences between
the decoded and ground-truth spectrograms. A reference loss
$\mathcal{L}_{\mathrm{ref}}$ penalises departures of the decoded speech
parameters from the reference parameters produced by the pre-trained
auto-encoder. A STOI-derived intelligibility loss complements these two.
The reference loss weights are: voice weight $\lambda_\alpha = 1.8$,
loudness $\lambda_L = 1.5$, pitch $\lambda_{f_0} = 0.4$; formant
frequencies $\lambda_{f_1}=3$, $\lambda_{f_2}=1.8$, $\lambda_{f_3}=1.2$,
$\lambda_{f_4}=0.9$, $\lambda_{f_5}=0.6$, $\lambda_{f_6}=0.3$; formant
amplitudes $\lambda_{a_1}=4$, $\lambda_{a_2}=2.4$, $\lambda_{a_3}=1.2$,
$\lambda_{a_4}=0.9$, $\lambda_{a_5}=0.6$, $\lambda_{a_6}=0.3$; and
broadband filter parameters $\lambda_{f_u}=10$, $\lambda_{a_u}=4$,
$\lambda_{b_u}=4$. The three outer losses are combined with weights
$\lambda_1=1.2$, $\lambda_2=0.1$, $\lambda_3=1$. All models are trained
with the Adam optimizer~\cite{kingma2015adam} using learning rate
$10^{-3}$, $\beta_1=0.9$, and $\beta_2=0.999$. Each participant's data is
split into $350$ training trials and $50$ test trials, and we additionally
report a word-level cross-validation in which $10$ unique words are held
out per fold over $5$ folds.

\subsection{Evaluation Metrics}

We evaluate decoding quality on held-out trials using two complementary
metrics. The Pearson correlation coefficient (PCC) between the decoded and
ground-truth spectrograms quantifies spectrotemporal fidelity and permits
direct comparison with prior ECoG-to-speech studies. The Short-Time
Objective Intelligibility Plus (STOI+) measure~\cite{taal2011algorithm}
applies a one-third octave band analysis with $15$ bands ranging from a
lowest centre frequency of $150\,\mathrm{Hz}$ to a highest of
approximately $4.3\,\mathrm{kHz}$, using $P=16$ segments, $M=2$, and
$N=30$ temporal frames, and captures perceptual intelligibility.
Statistical comparisons use the Wilcoxon signed-rank test for paired
comparisons (e.g., causal versus non-causal on the same participants) and
the Wilcoxon rank-sum test for unpaired comparisons (e.g., left versus
right hemisphere).

\subsection{Occlusion-Based Contribution Analysis}

To localise the cortical sources that drive each decoder, we performed an
occlusion analysis in which each electrode in turn was set to zero on the
testing data and the resulting reduction in decoding PCC was recorded.
Electrode-wise contributions were projected onto MNI anatomical
coordinates, averaged across participants, and normalised by local grid
density, yielding cortical maps that can be compared across decoders and
hemispheres.

\section{Results}

\subsection{Model architectures and causality}

Before we could ask whether causal ECoG decoding is feasible, we needed
a direct comparison of the principal deep-learning families currently in
use for neural speech decoding. We therefore trained all three decoder
backbones -- 3D ResNet, 3D Swin Transformer, and LSTM -- in both causal
and non-causal variants across all $N=48$ participants, and evaluated
them on the held-out trials using spectrogram PCC.

The ResNet decoder attained the highest mean PCC across participants
($\mathrm{PCC}=0.804$ non-causal, $\mathrm{PCC}=0.798$ causal), with the
Swin Transformer closely behind ($\mathrm{PCC}=0.793$ non-causal,
$\mathrm{PCC}=0.796$ causal) and the LSTM noticeably lower
($\mathrm{PCC}=0.757$ non-causal, $\mathrm{PCC}=0.713$ causal)
(Table~\ref{tab:architectures}, Figure~\ref{fig:pcc_models}). Paired
Wilcoxon signed-rank tests showed no significant difference between the
causal and non-causal ResNet ($p=0.1022$) or between the causal and
non-causal Swin ($p=0.2794$), while the causal LSTM was significantly
worse than the non-causal LSTM ($p=0.0007$). Direct comparison between
causal ResNet and causal Swin also failed to reach significance
($p=0.1413$), so the two convolutional and transformer backbones behave
comparably under causal constraints.

These findings matter because they clear the single largest obstacle to
deploying ECoG-to-speech decoders in real time. For the convolutional and
transformer families, strictly causal operation recovers essentially the
same spectrogram fidelity as the widely reported non-causal
configurations, removing the implicit reliance on post-utterance auditory
feedback. The LSTM exception, in contrast, suggests that the recurrent
family does exploit backward temporal context when permitted, and loses
material accuracy when it is not. We therefore focus subsequent analyses
on the causal ResNet and causal Swin decoders.

\begin{figure}[htbp]
\centering
\includegraphics[width=0.72\textwidth]{figures/fig_pcc_models.pdf}
\caption{Mean spectrogram Pearson correlation coefficient (PCC) across
$N=48$ participants for 3D ResNet, 3D Swin Transformer, and LSTM backbones,
evaluated in both non-causal and causal configurations. ResNet achieves the
highest PCC (non-causal $0.804$, causal $0.798$), closely followed by SWIN
($0.793$ and $0.796$). LSTM trails both (non-causal $0.757$, causal
$0.713$). Wilcoxon signed-rank tests comparing the causal and non-causal
variants give $p=0.1022$ for ResNet, $p=0.2794$ for SWIN, and $p=0.0007$ for
LSTM. Only LSTM shows a significant causal/non-causal gap.}
\label{fig:pcc_models}
\end{figure}

\begin{table}[htbp]
\centering
\caption{Decoding PCC across architectures and causality across $N=48$
participants. $p$-values are from Wilcoxon signed-rank tests comparing
causal and non-causal variants on the same participants.}
\label{tab:architectures}
\begin{tabular}{lccc}
\toprule
Architecture & Non-causal PCC & Causal PCC & Non-causal vs causal $p$ \\
\midrule
3D ResNet & $0.804$ & $0.798$ & $0.1022$ \\
3D Swin Transformer & $0.793$ & $0.796$ & $0.2794$ \\
LSTM & $0.757$ & $0.713$ & $0.0007$ \\
\bottomrule
\end{tabular}
\end{table}

\subsection{Speech parameter reconstruction}

Because the differentiable Speech Synthesizer makes every step of the
decoding chain interpretable, we next asked how well the ECoG Decoder
recovers each individual speech parameter. This analysis is important
because the spectrogram PCC summarises overall fidelity but can mask
systematic errors in specific acoustic features -- for example, a decoder
could produce a high-correlation spectrogram while systematically missing
pitch variation.

Using the causal ResNet decoder, we correlated the decoded and reference
time series for each speech parameter across participants. Pitch $f_0$
was recovered with a mean PCC of $0.907$, voice weight with $0.834$, the
first formant $f_1$ with $0.827$, the second formant $f_2$ with $0.895$,
and loudness with $0.545$ (Figure~\ref{fig:params}). The parameters most
tightly linked to intelligibility and participant-specific voice quality
-- pitch, voicing, and the first two formants -- are therefore recovered
most accurately, while loudness, which is known to be harder to read out
of cortical high-gamma activity, lags behind.

These parameter-level correlations explain why the reconstructed audio
preserves participant-specific voice quality rather than converging on a
generic vocoder voice. Because the Speech Synthesizer is fixed after
stage one, a well-reconstructed parameter stream directly translates into
a well-reconstructed waveform that retains the participant's pitch
contour and vowel formant structure. This behaviour contrasts with
kinematic latent spaces that tend to abstract away speaker
identity~\cite{anumanchipalli2019speech}.

\begin{figure}[htbp]
\centering
\includegraphics[width=0.72\textwidth]{figures/fig_params.pdf}
\caption{Reconstruction of decoded speech parameters using the causal ResNet
decoder. Mean Pearson correlation between decoded and reference time series
across participants: voice weight $0.834$, loudness $0.545$, pitch $f_0$
$0.907$, first formant $f_1$ $0.827$, second formant $f_2$ $0.895$.
Pitch, voice weight, and the first two formants -- the parameters most
important for intelligibility and participant-specific voice quality -- are
recovered most accurately.}
\label{fig:params}
\end{figure}

\begin{table}[htbp]
\centering
\caption{Decoded speech parameter reconstruction. Mean Pearson correlation
between decoded and reference time series across participants using the
causal ResNet decoder.}
\label{tab:params}
\begin{tabular}{lc}
\toprule
Speech parameter & Mean PCC across participants \\
\midrule
Voice weight ($\alpha$) & $0.834$ \\
Loudness ($L$) & $0.545$ \\
Pitch ($f_0$) & $0.907$ \\
First formant ($f_1$) & $0.827$ \\
Second formant ($f_2$) & $0.895$ \\
\bottomrule
\end{tabular}
\end{table}

\subsection{Left versus right hemisphere decoding}

Having established that causal ResNet and Swin decoders match their
non-causal counterparts, we next asked whether the framework generalises
from the conventionally studied left, language-dominant hemisphere to
the right hemisphere. This is clinically critical because many patients
who would benefit from a speech prosthesis have incurred left-hemisphere
damage, so the feasibility of right-hemisphere decoding directly
determines their eligibility for future devices.

We compared the $n=32$ left-hemisphere participants to the $n=16$
right-hemisphere participants using the causal ResNet and causal Swin
decoders. Right-hemisphere mean PCCs were $0.790$ (ResNet) and $0.781$
(Swin), closely tracking the whole-cohort causal means of $0.798$ and
$0.796$ (Table~\ref{tab:hemisphere}, Figure~\ref{fig:hemisphere}). Wilcoxon
rank-sum tests found no significant PCC difference between hemispheres
for either decoder (ResNet $p=0.312$; Swin $p=0.325$). The STOI+
intelligibility measure gave the same qualitative picture (ResNet
$p=0.108$; Swin $p=0.092$), and none of the four $p$-values crossed the
$\alpha=0.05$ threshold (dashed red line in
Figure~\ref{fig:hemisphere}).

The absence of a significant hemispheric difference, across two metrics
and two decoder families, constitutes the first systematic evidence that
modern deep decoders can extract speech from right-hemisphere cortex at
levels comparable to the left hemisphere. Clinically, this supports the
pursuit of right-hemisphere speech prostheses for the substantial
patient population with expressive aphasia secondary to dominant-hemisphere
injury.

\begin{figure}[htbp]
\centering
\includegraphics[width=0.95\textwidth]{figures/fig_hemisphere.pdf}
\caption{Right hemisphere decoding. Left: mean spectrogram PCC for causal
ResNet and causal Swin models across the full cohort (causal means $0.798$
and $0.796$ across $n=48$ participants) compared with the right-hemisphere
subset ($n=16$), which reaches PCC $0.790$ (ResNet) and $0.781$ (Swin).
Right: Wilcoxon rank-sum $p$-values for the left-vs-right comparison. PCC
p-values are $0.312$ (ResNet) and $0.325$ (Swin); STOI+ p-values are
$0.108$ (ResNet) and $0.092$ (Swin). None cross the $\alpha=0.05$ threshold
(dashed red line), consistent with the absence of significant hemispheric
differences.}
\label{fig:hemisphere}
\end{figure}

\begin{table}[htbp]
\centering
\caption{Right versus left hemisphere decoding with causal models.
$p$-values are from Wilcoxon rank-sum tests ($n=16$ right-hemisphere
participants versus $n=32$ left-hemisphere participants).}
\label{tab:hemisphere}
\begin{tabular}{lccc}
\toprule
Model & Right PCC & Left vs right $p$ (PCC) & STOI+ $p$ \\
\midrule
Causal ResNet & $0.790$ & $0.312$ & $0.108$ \\
Causal Swin & $0.781$ & $0.325$ & $0.092$ \\
\bottomrule
\end{tabular}
\end{table}

\subsection{Electrode density}

A further practical question for translation is whether clinical-grade
low-density ECoG grids, which are far more widely available than the
hybrid high-density arrays used in some research studies, can support
competitive decoding. We therefore compared the $n=5$ hybrid-density (HB)
participants with the $n=43$ low-density (LD) participants using the
causal ResNet and causal Swin decoders.

Mean PCC was similar across the two grid types, with ResNet performing
marginally better on the LD grids. To isolate the contribution of the
extra $1\,\mathrm{mm}$ micro-electrodes in the HB grids, we refit the
decoder using only the LD subset of electrodes within each HB
participant; for four of the five HB participants the PCC did not change
significantly, indicating that the additional micro-electrodes contribute
little under our current decoding framework. Across the full cohort,
$48\%$ of participants reached a decoding PCC between $0.80$ and $0.90$,
including participants with both grid types.

The practical implication is that the framework does not require
specialised high-density hardware to reach research-grade decoding
accuracy. This both simplifies future translational deployments, which
can proceed with commonly available clinical grids, and broadens the
cohort of existing epilepsy-surgery participants whose data can be
re-analysed with this pipeline.

\subsection{Cortical contribution analysis}

To connect decoder performance to the underlying neuroanatomy, we turned
to the occlusion-based contribution maps described in Section~3.7. These
maps reveal which cortical sites the decoder depends on most heavily and
therefore clarify what kind of information each model exploits.

Non-causal decoders placed substantial weight on superior temporal gyrus
(STG), consistent with the hypothesis that they benefit from
post-articulatory auditory feedback signals tracked by auditory cortex.
Causal decoders, by contrast, relied more heavily on frontal regions
including inferior frontal gyrus and ventral sensorimotor cortex, which
encode articulatory planning and execution rather than auditory feedback.
Right- and left-hemisphere causal models showed similar ventral
sensorimotor contributions, echoing the quantitative PCC result that
right-hemisphere decoding is indistinguishable from left-hemisphere
decoding at the cohort level.

These anatomical patterns provide a mechanistic interpretation of the
causal--non-causal PCC equivalence reported above. Causal decoders do
not simply reproduce non-causal decoders with a shifted time window;
they shift their representational burden from auditory to frontal cortex
and, in doing so, realise genuinely real-time-compatible speech
reconstruction.

\section{Discussion}

Our central finding is that an open, causal, deep-learning ECoG-to-speech
pipeline, built around a differentiable Speech Synthesizer, can achieve
spectrogram reconstruction fidelity equal to the best non-causal systems
and generalise across hemispheres and grid types. With the causal ResNet
backbone reaching a mean spectrogram PCC of $0.798$ across $N=48$
participants, and $48\%$ of those participants falling in the $0.80$--$0.90$
range, the framework meets or exceeds the accuracy of published
ECoG-to-speech systems while eliminating reliance on non-causal operations
that implicitly exploit auditory feedback. The two-stage training scheme,
in which a speech-only auto-encoder supplies reference speech parameters
to supervise the ECoG Decoder, is the key enabler: it turns the
otherwise data-starved task of ECoG-to-speech mapping into a much easier
cortex-to-parameter regression, while the differentiable synthesizer
still allows spectral loss to flow back end-to-end.

The results sharpen the understanding of the causal--non-causal gap
reported in earlier work~\cite{anumanchipalli2019speech,akbari2019towards,moses2021neuroprosthesis,metzger2023highperformance,silva2024bilingual}.
The near-identical ResNet and Swin results in causal and non-causal modes
($p=0.1022$ and $p=0.2794$) indicate that the gap closes completely for
convolutional and transformer architectures when those models have
adequate capacity to exploit frontal articulatory signals. The persistent
LSTM gap ($p=0.0007$) is consistent with the recurrent family's strong
sensitivity to backward temporal context, and suggests that the
increasingly common use of convolutional or transformer backbones in
modern speech BCIs~\cite{willett2023highperformance,metzger2023highperformance}
is not only a pragmatic choice but also a principled one from the
perspective of causality.

The parameter-level analysis provides a complementary lens on what the
decoder has learned. High PCCs for pitch ($0.907$), the second formant
($0.895$), voice weight ($0.834$), and the first formant ($0.827$) show
that the decoder recovers exactly those acoustic features most important
for intelligibility and speaker identity, while the lower loudness PCC
($0.545$) is consistent with the well-known difficulty of reading off
absolute amplitude from cortical high-gamma activity. Because the
synthesizer operates on these interpretable parameters rather than on an
articulatory kinematic space, reconstructed speech preserves
participant-specific voice quality, a property not delivered by kinematic
decoders~\cite{anumanchipalli2019speech}. The low-dimensional parameter
space also explains why the system trains reliably on the $400$ trials
available per participant, a regime in which richer latent spaces
typically over-fit.

The hemispheric comparison extends the reach of deep ECoG-to-speech
decoding beyond the language-dominant hemisphere. Right-hemisphere PCCs
of $0.790$ (ResNet) and $0.781$ (Swin), together with non-significant
left-versus-right tests on both PCC (ResNet $p=0.312$, Swin $p=0.325$)
and STOI+ (ResNet $p=0.108$, Swin $p=0.092$), indicate that the right
hemisphere carries decodable correlates of overt speech that are
quantitatively on par with those of the left hemisphere under this
framework. Clinically, this is the first systematic evidence that
supports considering right-hemisphere prostheses for the substantial
patient population with expressive aphasia following dominant-hemisphere
stroke or injury, complementing motor prostheses for other
effectors~\cite{moses2021neuroprosthesis,metzger2023highperformance,lorach2023walking}.
The occlusion analysis adds mechanistic weight to this conclusion by
showing that causal right- and left-hemisphere decoders recruit similar
ventral sensorimotor cortex, suggesting a shared articulatory basis.

Our work has several limitations that the framework itself does not
resolve. Training requires paired speech and ECoG recordings, which are
available only intra-operatively and thus constrain the cohort to
epilepsy-surgery participants producing overt speech; extending the
pipeline to attempted or imagined speech, as is needed for paralysed
users, will require careful adaptation and is not directly validated by
our data. The current evaluation focuses on word-level decoding over a
$50$-word set with $500\,\mathrm{ms}$ mean utterance duration, and
continuous sentence-level decoding at conversational speeds remains to
be demonstrated within this framework. Finally, the five hybrid-density
participants provide only preliminary evidence that low-density grids
match high-density grids; larger high-density cohorts would be needed to
rule out subtler differences.

Looking ahead, the open-source pipeline, the interchangeable ResNet,
Swin, and LSTM backbones, and the causal evaluation protocol together
provide a community reference point against which future decoders can be
benchmarked. The differentiable Speech Synthesizer is a natural target
for extension, for example by incorporating richer excitation models or
prosodic structure, and the parameter-space supervision strategy could
plausibly transfer to related low-data cortical decoding problems. Most
immediately, the demonstration that causal, bilaterally capable,
low-density ECoG decoding is achievable at current accuracy levels brings
the field one step closer to reproducible, real-time, natural-sounding
speech prostheses that work across a broad and clinically realistic
patient population.

\section{Conclusion}

We have introduced a two-stage neural speech decoding framework that
couples interchangeable deep ECoG Decoders to a fully differentiable
Speech Synthesizer with reference-parameter supervision. Across $N=48$
participants, the framework delivers causal spectrogram reconstruction at
a mean PCC of $0.798$ (ResNet) and $0.796$ (Swin), matches non-causal
performance for both backbones, and generalises from left to right
hemisphere and from hybrid-density to low-density clinical grids. By
releasing the complete pipeline as open-source software, we aim to
provide a reproducible reference platform for the community and to
broaden the population of patients for whom ECoG-based speech prostheses
may ultimately become feasible.

\bibliographystyle{unsrt}
\bibliography{references}

\end{document}
